\section{Quantitative guide to the regimes}
\label{app:regime-table}

Table~\ref{tab:all-regimes} records the cases of the proof in their logical
order. All discrepancy and eventual-absence conclusions concern the
retained sides of the stabilized counterexample subsequence, after the
total $o(N)$ label removals specified in the proof, and hold eventually
along that subsequence. Width always means $\log(N\max\mu)$, using
the original host-side size $N$. Thus a normalized atom bound $N\max\mu\le n^D$ gives width
$D\log n$. Defect means density of the other color. Availability means
that, for some fixed $\kappa>0$, a witness can be chosen after every
removal of at most $\kappa N$ labels from each side; the witness may
depend on the removed sets. See
Section~\ref{sec:general-conventions-and-probabilistic-tools}.
Every constant and exponent is fixed before $n\to\infty$, with the
local meaning assigned in its cited section. In particular, parameters
with the same name in different rows need not have the same value.

For the atom bounds below, $p_b$ denotes an odd probability row and $p_v$
an even probability row. A tuple posterior $\Pi_v$ is a law on $k$
first-side labels, so $N^k\max\Pi_v$ is its density bound relative to
the uniform product law. This differs from the one-label bound
$N\max p_v$. The limiting exponents $H^\dagger,H_L^\dagger$ are those
of Definition~\ref{def:limiting-bias-exponents}; they delimit the
available-bias ranges at power and linear width budgets.

The last six rows use initial discrepancy~\eqref{eq:source-2}, deep
discrepancy~\eqref{eq:source-9} in both orientations, and the eventual
cluster-witness absence of Corollary~\ref{cor:eventual-cluster-absence}.
Here $g,q$ are the residual bias and cluster scales of
Definition~\ref{def:residual-scales}; $M_i$ is the size of either side of
patch $i$, $\ell_i$ its prefix length, $h_i$ its internal dimension,
$d_i$ its bin size, and $s_i$ its gain parameter. The cleaned first-side
degrees are measured by
$D_{\pi_i}(x)=\sum_y\pi_i(y)\mathbf1_G(x,y)$, where $\pi_i$ is
uniform on $Y_i$ in the bounded and direct branches and is the comparison
law of Section~\ref{sec:internal-solver-and-profiles-for-these-tilings}
in the cluster branches. Every nonbounded row is reached only after
$\max(g,q)\ge M_1Q_0$. For such a patch the common allocation budget is
\[
 \ell_i,\ \log(N/M_i)\le s_i/(1000u),
\]
where $u$ is the fixed even moment order from
Section~\ref{sec:the-two-deep-orientations-some-estimates}. The word
\emph{bounded} names the first extraction branch; it does not assert that
the scales in every other branch tend to infinity.

\begingroup
\small
\setlength{\tabcolsep}{4pt}
\renewcommand{\arraystretch}{1.18}
\begin{longtable}{@{}>{\raggedright\arraybackslash}p{.255\linewidth}>{\raggedright\arraybackslash}p{.355\linewidth}>{\raggedright\arraybackslash}p{.345\linewidth}@{}}
\caption{All regimes: hypotheses, quantitative costs, and the interface
used by the next stage. In the first reductions an available witness
produces a cube; its exclusion supplies the stated discrepancy or
structural conclusion.}\label{tab:all-regimes}\\
\toprule
Regime and input & Width, density, and loss budget & Output and location\\
\midrule
\endfirsthead
\multicolumn{3}{l}{\textit{Table~\thetable\ continued}}\\
\toprule
Regime and input & Width, density, and loss budget & Output and location\\
\midrule
\endhead
\midrule
\multicolumn{3}{r}{\textit{Continued on the next page}}\\
\endfoot
\bottomrule
\endlastfoot

Counterexample sequence &
$n\to\infty$, $C_n=N/2^n\to\infty$, $N\le n2^n$. No rate is imposed
on $C_n$. &
The bipartite contradiction setting;
Lemma~\ref{lem:counterexample-sequence}.\\

Power-width bias versus purity &
Widths $(n^\beta,n^\gamma)$, $0<\beta\le\gamma<1$; an available
patch with absolute bias $\ge n^{-h}$, and no pair of laws on the
current sets with defect $\le e^{-n^h}$ in either color at the
respective doubled widths $(2n^\beta,2n^\gamma)$.
Here $\omega=\min(\beta,1-\gamma)/1000$, $h=\omega/10^6$.
Filter widths fit the plateau margin;
$N\max p_v\le\exp((\log2-c n^{-h})n)$. &
The two hypotheses together give a cube. The density gain pays for local
repetitions and predictive tests;
Lemma~\ref{lem:bias-purity}.\\

Constant broad side &
Balanced first laws of width $n^\gamma$, $0<\gamma<1$; at least $\chi N$
columns of degree $\ge.95$ for each tag. Given the parents, the stream
law relative to $R[u]$ is $\le e^{a_0u}$, $a_0=1/3$.
The posterior satisfies $N^k\max\Pi_v\le e^{\tau_0kn}$, and
$N\max p_v\le O(e^{\tau_1n})$, with
$a_0<\tau_0<\tau_1<\log2$. &
A cube in that color. The remaining gap below $\log2$ pays for geometric
overlap; Lemma~\ref{lem:constant-broad-high-neighborhood}.\\

Polynomial broad-side concentration and polynomial defect &
Balanced laws of widths at most $n^\gamma$ and $D_*\log n$, $0<\gamma<1$;
every broad-side degree is $\ge1-n^{-p_0}$, for fixed $p_0>0$.
After marginalizing the tag, $N^k\max\Pi_v\le e^{.4kn}$;
$N\max p_v\le O(e^{.55n})$. The constant $D_*>0$ is universal. &
A cube in one of the colors;
Lemma~\ref{lem:small-polynomial-broad-side}. This excludes the sparse
rectangles tested in~\eqref{eq:source-1}.\\

Small-grid purity &
Available patches have both widths $\le n^{d/2}$ and every second-support
degree $\ge1-e^{-n^p}$, $p>0$. Every law pair of widths
$(D_0\log n,n^{1/4})$ has both color densities $>n^{-c}$, where
$0<D_0<\min(.1,D_*/2)$.
Use $d<D_0/1000$, $b=d/4$, $c=b/20$;
$\log(N\max p_b)\le n^{d/2}+2cn^d\log n+O(\log n)$ and
$\log(N\max p_v)\le .04n^{4d}+o(n^{4d})$. &
Purity gives a cube. Together with the power-width alternative this yields
discrepancy $n^{-\eta_0}$ at both widths $n^{\eta_0}$, for a fixed
$\eta_0>0$; Lemma~\ref{lem:small-grid-purity} and
Corollary~\ref{cor:initial-discrepancy}.\\

Asymmetric purity &
Under initial discrepancy~\eqref{eq:source-2}, a finite balanced
mixture has widths $(n^\gamma,n^\beta)$, $\gamma<1$, $0<\beta<\tau/4$, where $\tau=\min(\eta_0/2,.04)/4$;
every second-support degree is $\ge1-e^{-n^p}$ for fixed $p>0$.
With $s=\lceil n^\tau\rceil$, $N\max p_b\le e^{.02s\log n}$
and $N\max p_v\le e^{.1n}$.
No relation $p>\tau$ is required. &
Purity gives a cube. Consequently, at every admissible fixed rational
width pair, all weighted densities differ from $1/2$ by at most a fixed
negative power of $n$, in both orientations;
Lemma~\ref{lem:asymmetric-purity} and
Corollary~\ref{cor:asymmetric-discrepancy}.\\

Intermediate sublinear bias: some $H^\dagger<1$ &
Shallow widths $(n^{x_s},n^{y_s})$ and deep widths
$(n^{x_d},n^{y_d})$, $y_s<1-\sigma<y_d$.
Available bias $2a_*=n^{-h_+}$; deep error $b_*=n^{-h_-}$,
$0<h_-<h_+<1$, $h_+-h_-<\chi/10$.
The gain $cna_*$ pays for first-side width and
$O(n^\sigma\log n)$ local entropy. &
A cube, using a core of size $\le n^\chi$ and an outside-core read bound
$n^\sigma+3$; Proposition~\ref{prop:intermediate-bias-exclusion}.
This excludes the whole range $H^\dagger<1$, including zero.\\

Intermediate linear bias: both $H^\dagger\ge1$, some
$0<H_L^\dagger<1$ &
Shallow widths $(n^{x_s},\alpha_s n)$ and deep widths
$(n^{x_d},\alpha_d n)$, $100\alpha_s<\alpha_d$.
Use $a_*=\tfrac12n^{-h_+}$, $b_*=n^{-h_-}$ with parameters chosen
afresh in the linear case; $\sigma<\chi/10$ and broad-test error
$o(a_*)$. The read-tail exponent is
$na_*^2/((n^\sigma+3)b_*^2)$. &
A cube. The surplus survives the external-coordinate costs;
Proposition~\ref{prop:intermediate-bias-exclusion}.
The two intermediate cases use different parameter orders.\\

Full-dimensional clusters &
Assume initial discrepancy in both orientations and available one-color
patches. The first law and aggregate second law have widths
$\le n^\delta$; individual cluster atoms are $\le e^{-n^\zeta}$.
Codegrees, including repeated labels, are $\ge1/4+n^{-\delta}$.
Here $\delta<\min(\eta_0,\zeta,1)/2000$ and
$k(T+m)=O(n^{500\delta})$ fits both discrepancy and atom budgets.
The even bound is $N\max p_v\le e^{(\log2-.01n^{-\delta})n}$. &
A cube, hence eventual absence for every fixed admissible rational
choice. Used in the linear jump and again in extraction;
Proposition~\ref{prop:full-dimensional-cluster-exclusion} and
Corollary~\ref{cor:eventual-cluster-absence}.\\

Linear-budget jump: both $H^\dagger\ge1$, some $H_L^\dagger=0$ &
Available widths $(n^{.01},.011n)$ with column surplus $2n^{-.01}$.
Reverse-orientation discrepancy $n^{-.95}$ at widths
$(n^{1-\upsilon/4},n^{x_0})$.
The internal dimension $n^{.1}$ gives gain $n^{.09}$, exceeding
$O(n^{.05})+n^{.03}$ losses. &
A cube; Proposition~\ref{prop:linear-budget-jump-exclusion}.
Thus both $H_L^\dagger\ge1$ remain.\\

Deep discrepancy &
Widths $(\alpha n,n^{x_*})$ in either order, for fixed
$0<\alpha,x_*<.01$; error $b_*=n^{-1+.04}$.
Every fixed smaller positive error slack is available with possibly
smaller budgets. The interaction tails use the setting of
Lemma~\ref{lem:heterogeneous-centered-moment}: the first law $\tau$
and second laws $\pi_l$ have widths $\le n^{x_*/4}$, with
$|D_{\pi_l}(x)-1/2|\le O(b_*)$ on $\operatorname{supp}\tau$.
For each fixed $u$, the one-free threshold $O_u(b_*)$ has tail
$e^{-\alpha n/3}$, and the two-free threshold $n^{-1-.03}$ has tail
$e^{-n^{.4}}$. &
Interaction moments, row trimming, and homogeneous clique peeling;
Corollary~\ref{cor:deep-discrepancy} and
Lemmas~\ref{lem:deep-interaction-tails}--\ref{lem:simultaneous-row-trimming}.
Interaction bounds allow different laws at different coordinates;
clique peeling uses a single common law.\\

Bounded extraction: $\max(g,q)<M_1Q_0$ &
One patch $M_i\asymp N$, $\ell_i=h_i=s_i=0$.
On cleaned supports, $|D_{\pi_i}(x)-1/2|\le K/n$ and the number
of labels with a large interaction is bounded by a fixed constant.
The $r\asymp A_0\log n$ late classes supply polynomial slack. &
Two surviving candidates per even role, with joint bounds that give Hall;
Proposition~\ref{prop:patch-extraction-allocation},
Lemma~\ref{lem:stable-patch-cleaning}, and
Sections~\ref{sec:low-modes-geometry-and-cell-states}--\ref{sec:late-dynamics-and-its-transfer-problem}.\\

Low direct: $g>M_1q$, $g\le K_B\log n$ &
$M_i\ge cNe^{-g^{a_B}}$, $h_i=0$, $s_i=10^{-3}g$.
On cleaned supports,
$1/2+g/(4n)\le D_{\pi_i}(x)\le1/2+4g/n$.
Prefix and width losses are $\le s_i/(1000u)$; late classes pay for
the remaining list losses. &
Two surviving candidates per even role, with joint bounds that give Hall;
Proposition~\ref{prop:patch-extraction-allocation},
Lemma~\ref{lem:stable-patch-cleaning}, and
Lemma~\ref{lem:low-two-endpoint-hall}.\\

High direct: $g>M_1q$, $g>K_B\log n$ &
The same direct patch budgets and degree interval.
The homogeneous peel parameter is $\le e^{-200s_i}$, giving mass-failure
probability $\le n^{-R}$. The unnormalized row bound
$M_iZ_v(x)\le2^n e^{-200s_i}$ pays for repeated local neighborhoods. &
An odd injection and even weights $Z_v$ supported on common neighbors,
with row masses $\ge1/2$ and column sums $\le1/2$. Normalizing gives Hall;
Lemma~\ref{lem:high-direct-assignment}.\\

Low cluster: $g\le M_1q$, $Q_0\le q\le(\log n)^{c_q}$ &
$M_i\ge cNe^{-q^{a_C}}$, $d_i=\lfloor e^{q/2}\rfloor$,
$q^{M_{\rm lo}}\le h_i<2q^{M_{\rm lo}}$,
$s_i=ah_i/10^6$.
Within-bin codegree is $\ge1/4+a$; on cleaned supports,
$|D_{\pi_i}(x)-1/2|\le10q^{C_b}/n$, with
$10q^{C_b}<s_i/(100u)$. Also $h_i\le(\log n)^{.1}$. &
Exact group and label marginals preserve the mean estimates.
The later assignment gives two candidates per even role and Hall;
Proposition~\ref{prop:internal-slice-solver},
Propositions~\ref{prop:calibrated-group-bins}
and~\ref{prop:calibrated-role-labels},
Lemma~\ref{lem:low-two-endpoint-hall}.
The scales $q,h_i,d_i$ may stay fixed.\\

High cluster with small bins: $g\le M_1q$,
$(\log n)^{c_q}<q\le(\log n)^2$ &
The same cluster width, gain, degree, and codegree bounds, with
$q^{M_{\rm hi}}\le h_i<2q^{M_{\rm hi}}$; $d_i$ is the smaller of
$\lfloor e^{q/2}\rfloor$ and $\lfloor e^{\sqrt{\log n}}\rfloor$.
Here $h_i>(\log n)^5$ and $h_i=o(d_i^{.025})$.
Exact singleton and calibrated joint estimates control repeated bins. &
An odd injection and even weights $Z_v$ with row masses $\ge1/2$
and column sums $\le1/2$. Normalizing the rows gives Hall;
Propositions~\ref{prop:high-cluster-sampling}
and~\ref{prop:high-cluster-loads}.\\

High cluster with large bins: $g\le M_1q$, $q>(\log n)^2$ &
The same cluster budgets, now $d_i=\lfloor e^{q/2}\rfloor$ and
$q^{M_{\rm hi}}\le h_i<2q^{M_{\rm hi}}$.
With $k_i=\lceil h_i^{3\omega}\rceil$,
$T_i=\lceil h_i^\omega\rceil$, all cluster modes have
$k_iT_i\ll q^{a_C/2},\log d_i$.
The conditional in-bin label law satisfies
$\max_y U_g^W(D)(y)\le2e^{1.5k_iT_i}/d_i$, a superpolynomially
small bound. &
An odd injection and even weights $Z_v$ with row masses $\ge1/2$
and column sums $\le1/2$. Normalizing the rows gives Hall;
Propositions~\ref{prop:high-cluster-sampling}
and~\ref{prop:high-cluster-loads}.\\
\end{longtable}
\endgroup

\paragraph{Exhaustion and boundary conventions.}
The intermediate-bias proposition first excludes $H^\dagger<1$ in either
orientation. Once both $H^\dagger\ge1$, each $H_L^\dagger$ belongs to
$\{0\}$, $(0,1)$, or $[1,+\infty]$. The first two possibilities are
excluded separately, and the third gives deep discrepancy by taking
strictly smaller error exponents. No conclusion at the threshold defining
an infimum is needed.

During extraction the bounded case is tested first. Outside it,
$g\le M_1q$ implies $q\ge Q_0$, so the cluster branch is defined.
Equality at $g=M_1q$ belongs to the cluster branch, equality at
$g=K_B\log n$ to low direct, and equality at either cluster threshold
to the lower range. The fixed choice of $Q_0$ makes estimates in its
positive powers uniformly sufficient; it never licenses replacing a
fixed-scale error by an error tending to zero with $n$.
